\section{Array Lanjutan: map, filter, reduce, dan Metode Fungsional}

Metode fungsional array menghasilkan array atau nilai baru \textbf{tanpa} mengubah array asli (\textit{immutable pattern}). Ini menjadikan kode lebih mudah dipahami dan di-debug. Tiga metode terpenting: \code{map()} mentransformasi setiap elemen, \code{filter()} menyaring elemen berdasarkan kondisi, dan \code{reduce()} mengakumulasi semua elemen menjadi satu nilai \cite{jsInfo}.

\begin{table}[htbp]
\centering
\begin{tabular}{|P{2.5cm}|P{4.5cm}|P{5cm}|}
\hline
\textbf{Metode} & \textbf{Fungsi} & \textbf{Contoh} \\
\hline
\code{.map(fn)} & Transformasi setiap elemen & \code{[1,2,3].map(x => x*2)} $\rightarrow$ \code{[2,4,6]} \\
\hline
\code{.filter(fn)} & Saring berdasarkan kondisi & \code{[1,2,3].filter(x => x>1)} $\rightarrow$ \code{[2,3]} \\
\hline
\code{.reduce(fn, init)} & Akumulasi ke satu nilai & \code{[1,2,3].reduce((s,x)=>s+x,\allowbreak 0)} $\rightarrow$ 6 \\
\hline
\code{.find(fn)} & Elemen pertama yang cocok & \code{arr.find(x => x.id===5)} \\
\hline
\code{.some(fn)} & Minimal 1 cocok? (boolean) & \code{arr.some(x => x > 10)} \\
\hline
\code{.every(fn)} & Semua cocok? (boolean) & \code{arr.every(x => x > 0)} \\
\hline
\code{.sort(fn)} & Urutkan (mutasi!) & \code{arr.sort((a,b) => a-b)} \\
\hline
\code{.flat(depth)} & Ratakan array bersarang & \code{[[1],[2,3]].flat()} $\rightarrow$ \code{[1,2,3]} \\
\hline
\end{tabular}
\caption{Metode array fungsional dan contoh penggunaannya}
\end{table}

\begin{webcode}[language=JavaScript, caption={Chaining metode array}]
const products = [
  { name: "Laptop", price: 12000000, stock: 5 },
  { name: "Mouse",  price: 150000,   stock: 0 },
  { name: "Monitor",price: 3500000,  stock: 8 }
];

const result = products
  .filter(p => p.stock > 0)           // hanya yang tersedia
  .map(p => ({ ...p, tax: p.price * 0.1 })) // tambah pajak
  .sort((a, b) => a.price - b.price); // urutkan harga

const total = result.reduce((sum, p) => sum + p.price, 0);
\end{webcode}

\begin{catatan}
\textbf{sort() Bermutasi!} Berbeda dari \code{map}/\code{filter}/\code{reduce}, \code{sort()} mengubah array asli. Untuk menghindari mutasi, salin dulu: \code{[...arr].sort(fn)}. Tanpa fungsi komparator, \code{sort()} mengurutkan secara alfabetis (bahkan angka: \code{[10,2,1].sort()} $\rightarrow$ \code{[1,10,2]}). Selalu berikan komparator untuk angka: \code{(a,b) => a-b} \cite{mdnJsGuide}.
\end{catatan}

